
de $Y$
dans $X$, supposé partout de rang $d$. Dans loc. cit., cette formule résulte d'un calcul très simple de faisceaux $\underline{\Tor}$. Grâce à ce qu'on a vu dans 4.2, la formule précédente implique la même formule dans $A(Y)$, mais seulement \uline{modulo torsion}. Il est assez extraordinaire que la formule (4.6) dans $A(Y)$ lui-même ne soit pas prouvée à l'heure actuelle\SGAAuthorFootnote{*}{cf. note de la page 23.\newline (**)}, ni même le cas particulier obtenu en faisant $y=1$, donnant la self-intersection de $Y$:
\begin{equation}\tag{4.7}
 f^*(\cl(Y))=c_d(\underline{N}),
\end{equation}
même si $k$ est le corps des complexes !

On se heurte à un problème analogue lorsqu'on tente de calculer l'anneau de Chow de la variété éclatée $X'$ de $X$ le long de $Y$, en calquant la méthode suivie pour le calcul de $\operatorname{Gr}^\bullet(X')_\Q$ [Exp. VII 4.8], et qui marche également (sans tensoriser par $\Q$) pour le calcul de $\Gtop^\bullet(X')$. La formule-clef pour cette détermination (avec les notations de loc. cit.) est
\begin{equation}\tag{4.8}
 f^*(i_*(y))=j_*\bigl(g^*(y)c_{d-1}(\check{\underline F})\bigr),
\end{equation}
formule établie à l'heure actuelle dans $\Gtop^\bullet(X')$ (pour $y\in\Gtop^\bullet(Y)$) grâce à un calcul de $\underline{\Tor}$ (Exp. VII 3.4), et dont la validité dans $A(X')$ (pour $y\in A(Y)$) resterait à établir.

Signalons que même les analogues $\ell$-adiques des formules (4.6), (4.7), (4.8) ne sont pas prouvés à l'heure actuelle; lorsque $y$ est la classe de cohomologie d'un cycle algébrique, (4.6) et (4.8) sont vraies en tous cas mod. torsion.

\textbf{4.5.} Dans cette section et les suivantes, nous allons donner des exemples où dans le diagramme (4.1) l'homomorphisme $j$ respectivement l'homomorphisme $\phi$ n'est pas un isomorphisme. Il revient évidemment au même, pour $X$ donné, de dire que
\[
 j:\operatorname{Gr}^\bullet(X)\longrightarrow \Gtop^\bullet(X)
\]
est un isomorphisme, ou un monomorphisme, ou un épimorphisme. D'ailleurs l'image de $j$ est le sous-anneau de $\Gtop^\bullet(X)$ engendré par les classes de Chern des faisceaux localement libres sur $X$. Dans le cas où $K^\bullet(X)$ est engendré comme algébre sur $\Z$ par les classes des faisceaux inversibles sur $X$, on voit tout de suite que le sous-anneau envisagé de $\Gtop^\bullet(X)$ est aussi le sous-anneau engendré par $\Gtop^1(X)$, la composante de degré $1$ de $\Gtop^\bullet(X)$.
